% Hafta 9 — Sanallaştırma: en güçlü, en pahalı
% Derleme: py -3.12 tools/latex/derle.py docs/week-9/assets/h09-06-sanallastirma.tex
\documentclass[tikz,border=8pt]{standalone}
\usepackage{cen429-cizim}
\begin{document}
\begin{tikzpicture}[node distance=10mm and 10mm]
  \node[baslik] (b) at (0,0) {Sanallaştırma: en güçlü, en pahalı};
  \node[altbaslik, text width=170mm] at (b.south west) {analizci artık sizin uydurduğunuz makineyi çözmek zorunda};

  \node[kutu=36mm, below=10mm of b.south west, anchor=north west] (k1) {\kb{Orijinal fonksiyon}\\makine kodu};
  \node[koyukutu=36mm, right=of k1] (k2) {\kb{Bytecode'a çevir}\\özel komut kümesi};
  \node[koyukutu=36mm, right=of k2] (k3) {\kb{VM (yorumlayıcı)}\\bytecode'u çalıştırır};
  \node[iyikutu=36mm, right=of k3] (k4) {\kb{Sonuç aynı}\\ama makine kodu yok};
  \draw[ok, draw=soluk] (k1.east) -- (k2.west);
  \draw[ok, draw=soluk] (k2.east) -- (k3.west);
  \draw[ok, draw=soluk] (k3.east) -- (k4.west);

  \node[kotudip, text width=170mm, below=10mm of k1.south -| k1, anchor=north west] {%
    Maliyet: onlarca kat yavaşlama + boyut $\rightarrow$ yalnız küçük ve kritik fonksiyonlara.};
\end{tikzpicture}
\end{document}
